% paper_rtsc_tau_supplement.tex
% Supplemental Material for:
% Room-Temperature Superconductivity as a Quantum Information Problem
% Author: Sheng-Kai Huang
% Compile: pdflatex paper_rtsc_tau_supplement.tex && pdflatex paper_rtsc_tau_supplement.tex

\documentclass[prb,twocolumn,superscriptaddress,longbibliography,floatfix]{revtex4-2}

\usepackage{amsmath}
\usepackage{amssymb}
\usepackage{amsfonts}
\usepackage{amsthm}
\usepackage{bm}
\usepackage{hyperref}
\usepackage{booktabs}
\usepackage{graphicx}
\usepackage{listings}
\usepackage{xcolor}

\newtheorem{theorem}{Theorem}
\newtheorem{corollary}[theorem]{Corollary}
\newtheorem{lemma}[theorem]{Lemma}
\newtheorem{definition}[theorem]{Definition}
\newtheorem*{proposition}{Proposition}

% Shared macros (Paper 1 conventions)
\newcommand{\kB}{k_{\mathrm{B}}}
\newcommand{\Tr}{\mathrm{Tr}}
\newcommand{\Petz}{\mathcal{R}}
\newcommand{\chan}{\mathcal{N}}
\newcommand{\ket}[1]{|#1\rangle}
\newcommand{\bra}[1]{\langle#1|}
\newcommand{\braket}[2]{\langle#1|#2\rangle}
\newcommand{\op}[2]{|#1\rangle\langle#2|}
\newcommand{\abs}[1]{\left|#1\right|}
\newcommand{\tauasym}{\tau}
% Paper-specific macros
\newcommand{\taueff}{\tauasym_{\mathrm{eff}}}
\newcommand{\taudec}{\tauasym_{\mathrm{ph}}}
\newcommand{\Icoop}{I_{\mathrm{Cooper}}}
\newcommand{\Ipair}{I_{\mathrm{pair}}}
\newcommand{\Rtau}{R_{\tauasym}}
\newcommand{\Tc}{T_{\mathrm{c}}}
\newcommand{\oD}{\omega_{\mathrm{D}}}
\newcommand{\EF}{E_{\mathrm{F}}}
\newcommand{\Ek}{E_{\bm{k}}}
\newcommand{\xik}{\xi_{\bm{k}}}
\newcommand{\vF}{v_{\mathrm{F}}}
\newcommand{\ThetaD}{\Theta_{\mathrm{D}}}

% Listings style for Python code
\lstset{
  language=Python,
  basicstyle=\ttfamily\footnotesize,
  keywordstyle=\color{blue},
  commentstyle=\color{gray},
  stringstyle=\color{red!60!black},
  numbers=left,
  numberstyle=\tiny\color{gray},
  numbersep=5pt,
  breaklines=true,
  frame=single,
  captionpos=b
}

\begin{document}

\title{Supplemental Material for ``Room-Temperature Superconductivity as a
Quantum Information Problem: The $R_\tau$ Screening Criterion and
Material Predictions''}

\author{Sheng-Kai Huang}
\email{akai@fawstudio.com}
\affiliation{Independent Researcher}

\date{\today}

\maketitle

This Supplemental Material provides detailed derivations
(Secs.~S1--S5), material-specific calculations (Sec.~S6),
experimental protocols (Sec.~S7), and the $\Rtau$ calculator
code (Sec.~S8).

%=======================================================================
\section{S1: Derivation of $\taudec(T)$ from the electron-phonon
Hamiltonian}
\label{sec:S1}
%=======================================================================

\subsection{S1.1: Setup}

The electron-phonon Hamiltonian is
\begin{equation}\label{eq:S1_Hep}
  H = H_{\mathrm{el}} + H_{\mathrm{ph}} + H_{\mathrm{ep}}\,,
\end{equation}
where
\begin{align}
  H_{\mathrm{el}} &= \sum_{\bm{k}\sigma} \epsilon_{\bm{k}}\,
  c_{\bm{k}\sigma}^\dagger c_{\bm{k}\sigma}\,, \\
  H_{\mathrm{ph}} &= \sum_{\bm{q}\lambda} \hbar\omega_{\bm{q}\lambda}\,
  a_{\bm{q}\lambda}^\dagger a_{\bm{q}\lambda}\,, \\
  H_{\mathrm{ep}} &= \sum_{\bm{k}\bm{q}\lambda\sigma}
  g_{\bm{k},\bm{q}\lambda}\,
  c_{\bm{k}+\bm{q},\sigma}^\dagger c_{\bm{k}\sigma}\,
  (a_{\bm{q}\lambda} + a_{-\bm{q}\lambda}^\dagger)\,.
\end{align}
The electron-phonon coupling matrix element $g_{\bm{k},\bm{q}\lambda}$
encodes the scattering amplitude for an electron at $\bm{k}$
to scatter to $\bm{k}+\bm{q}$ by emitting or absorbing a phonon
in branch $\lambda$.

\subsection{S1.2: Scattering rate from Fermi golden rule}

The transition rate for scattering $\bm{k} \to \bm{k}' = \bm{k}+\bm{q}$
via phonon $(\bm{q},\lambda)$ is
\begin{align}\label{eq:S1_rate}
  W_{\bm{k}\to\bm{k}'}^{(\pm)} &= \frac{2\pi}{\hbar}\,
  |g_{\bm{k},\bm{q}\lambda}|^2\,
  \bigl(n_{\bm{q}\lambda} + \tfrac{1}{2} \pm \tfrac{1}{2}\bigr)
  \nonumber\\
  &\quad \times \delta(\epsilon_{\bm{k}'} - \epsilon_{\bm{k}}
  \mp \hbar\omega_{\bm{q}\lambda})\,,
\end{align}
where $+$ denotes phonon emission and $-$ phonon absorption,
and $n_{\bm{q}\lambda} = (e^{\hbar\omega_{\bm{q}\lambda}/\kB T} - 1)^{-1}$.

\subsection{S1.3: Entropy production per scattering event}

The entropy production per scattering event follows from
the detailed balance ratio.  For emission ($\bm{k} \to \bm{k}'$
with phonon created):
\begin{equation}
  \frac{W_{\mathrm{em}}}{W_{\mathrm{abs}}}
  = \frac{n_{\bm{q}\lambda} + 1}{n_{\bm{q}\lambda}}
  = e^{\hbar\omega_{\bm{q}\lambda}/\kB T}\,.
\end{equation}
The entropy produced is
\begin{equation}\label{eq:S1_sigma_event}
  \Sigma_{\mathrm{event}} = \ln\frac{W_{\mathrm{em}}}{W_{\mathrm{abs}}}
  = \frac{\hbar\omega_{\bm{q}\lambda}}{\kB T}\,.
\end{equation}
This has units of nats and is dimensionally correct:
$[\hbar\omega/\kB T] = [\mathrm{energy}/\mathrm{energy}]
= \text{dimensionless}$.

\subsection{S1.4: Thermal average over phonon spectrum}

The Eliashberg spectral function is defined as
\begin{equation}
  \alpha^2 F(\omega) = N(0) \sum_{\bm{q}\lambda}
  |g_{\bm{q}\lambda}|^2\, \delta(\omega - \omega_{\bm{q}\lambda})\,,
\end{equation}
and the electron-phonon coupling constant is
\begin{equation}
  \lambda = 2\int_0^{\omega_{\max}} \frac{\alpha^2 F(\omega)}{\omega}\, d\omega\,.
\end{equation}

The thermally averaged entropy production per mean free path
involves averaging $\Sigma_{\mathrm{event}}$ over the phonon
spectrum weighted by the scattering probability:
\begin{align}\label{eq:S1_Sigma_avg}
  \overline{\Sigma}_{\mathrm{ep}}(T)
  &= \frac{\int_0^{\omega_{\max}} \alpha^2 F(\omega)\,
  \frac{\hbar\omega}{\kB T}\, \coth\!\left(\frac{\hbar\omega}{2\kB T}\right)
  d\omega}
  {\int_0^{\omega_{\max}} \alpha^2 F(\omega)\,
  \coth\!\left(\frac{\hbar\omega}{2\kB T}\right) d\omega}\,.
\end{align}
In the Debye model, $\alpha^2 F(\omega) = (\lambda/2)\,
\omega/\oD$ for $\omega \leq \oD$ and zero otherwise.
Substituting and defining $x = \hbar\omega/\kB T$:
\begin{align}
  \overline{\Sigma}_{\mathrm{ep}}(T)
  &= \frac{\int_0^{x_D} x^2 \coth(x/2)\, dx}
  {\int_0^{x_D} x\, \coth(x/2)\, dx}\,,
  \quad x_D = \frac{\ThetaD}{T}\,.
\end{align}
However, this ratio is $O(1)$ and does not capture the overall
magnitude.  The total scattering entropy production per mean
free path (summing over all scattering events per $\ell$)
is obtained by multiplying by the average number of scattering
events per $\ell$, which is $\lambda\, \kB T/(\hbar\oD)$ for
$T > \ThetaD$.  The complete result is:
\begin{equation}\label{eq:S1_Sigma_full}
  \overline{\Sigma}_{\mathrm{ep}}(T)
  = \frac{2\pi\lambda\, \kB T}{\hbar\oD}\, \Phi\!\left(\frac{\ThetaD}{T}\right),
\end{equation}
where the correction factor is
\begin{equation}\label{eq:S1_Phi}
  \Phi(y) = \frac{2}{y^2}\int_0^{y} \frac{t^2\, dt}{e^t - 1}\,.
\end{equation}

\textbf{Limiting cases:}

\noindent\emph{High-$T$ limit} ($T \gg \ThetaD$, $y \to 0$):
$e^t - 1 \approx t$, so $\Phi(y) \to (2/y^2)\int_0^y t\, dt = 1$.
Thus
\begin{equation}
  \overline{\Sigma}_{\mathrm{ep}} \to \frac{2\pi\lambda\, \kB T}{\hbar\oD}
  = 2\pi\lambda\, \frac{T}{\ThetaD}\,.
\end{equation}

\noindent\emph{Low-$T$ limit} ($T \ll \ThetaD$, $y \to \infty$):
$\int_0^\infty t^2/(e^t-1)\, dt = 2\zeta(3) \approx 2.404$.
Thus
$\Phi(y) \to 2 \cdot 2.404 / y^2 = 4.808\, T^2/\ThetaD^2$, and
\begin{equation}
  \overline{\Sigma}_{\mathrm{ep}} \to
  \frac{2\pi\lambda \cdot 4.808\, \kB T^3}{\hbar\oD\, \ThetaD^2}
  = 9.62\pi\lambda\, \frac{T^3}{\ThetaD^3}\,,
\end{equation}
recovering the $T^3$ Bloch--Gr\"uneisen behavior.

\subsection{S1.5: From $\overline{\Sigma}$ to $\taudec$}

Using the Petz recovery bound $F \geq e^{-\Sigma/2}$
(Refs.~\cite{Junge2018,Fawzi2015} of the main text):
\begin{equation}\label{eq:S1_tau}
  \taudec(T) = 1 - e^{-\overline{\Sigma}_{\mathrm{ep}}(T)/2}\,.
\end{equation}
For small $\overline{\Sigma}$,
$\taudec \approx \overline{\Sigma}/2$.  For large
$\overline{\Sigma}$, $\taudec \to 1$.

\textbf{Numerical check} (Cu at $T = 300\,$K):
$\lambda = 0.13$, $\ThetaD = 343\,$K, $T/\ThetaD = 0.87$,
$\Phi(1.14) \approx 0.78$.
$\overline{\Sigma} = 2\pi(0.13)(0.87)(0.78) = 0.55$.
$\taudec = 1 - e^{-0.28} = 0.24$.
The experimental residual-resistivity-ratio-corrected value is
$\taudec \approx 0.16$--$0.25$, consistent.

%=======================================================================
\section{S2: Cooper pair entanglement calculation}
\label{sec:S2}
%=======================================================================

\subsection{S2.1: BCS ground state and Schmidt decomposition}

The BCS ground state factorizes over momentum modes:
\begin{equation}
  \ket{\mathrm{BCS}} = \prod_{\bm{k}}
  \bigl(u_{\bm{k}}\ket{0_{\bm{k}\uparrow} 0_{-\bm{k}\downarrow}}
  + v_{\bm{k}}\ket{1_{\bm{k}\uparrow} 1_{-\bm{k}\downarrow}}\bigr)\,.
\end{equation}
Each mode pair $(\bm{k}\!\uparrow, -\bm{k}\!\downarrow)$ is
already in Schmidt form with coefficients $|u_{\bm{k}}|$ and
$|v_{\bm{k}}|$.  The reduced density matrix for electron
$\bm{k}\!\uparrow$ is
\begin{equation}
  \rho_{\bm{k}\uparrow} = |u_{\bm{k}}|^2\, \op{0}{0}
  + |v_{\bm{k}}|^2\, \op{1}{1}\,,
\end{equation}
a diagonal state with eigenvalues $(|u_{\bm{k}}|^2, |v_{\bm{k}}|^2)$.

\subsection{S2.2: Entanglement entropy per mode}

The von Neumann entropy of $\rho_{\bm{k}\uparrow}$ gives the
entanglement entropy:
\begin{equation}\label{eq:S2_Sk}
  S_{\bm{k}} = H_{\mathrm{bin}}(p)\,,
  \quad p = |v_{\bm{k}}|^2 = \frac{1}{2}\left(1 - \frac{\xik}{\Ek}\right),
\end{equation}
where $H_{\mathrm{bin}}(p) = -p\ln p - (1-p)\ln(1-p)$ is the
binary entropy function (in nats).

Properties of $S_{\bm{k}}$:
\begin{itemize}
\item At the Fermi surface ($\xik = 0$): $p = 1/2$,
  $S_{\bm{k}} = \ln 2 \approx 0.693$ nats (maximally entangled).
\item Far from the Fermi surface ($|\xik| \gg \Delta$):
  $p \to 0$ or $p \to 1$, $S_{\bm{k}} \to 0$ (unentangled).
\item Width of the entangled region: $|\xik| \lesssim \Delta$.
\end{itemize}

\subsection{S2.3: Total Cooper pair entanglement}

Converting the sum over $\bm{k}$ to an energy integral using the
density of states $N(\xi) \approx N(0)$ (constant near $\EF$):
\begin{align}\label{eq:S2_Itotal}
  \Icoop &= N(0) \int_{-\hbar\oD}^{\hbar\oD}
  H_{\mathrm{bin}}\!\left(\frac{1}{2}\left[1-\frac{\xi}{\sqrt{\xi^2+\Delta^2}}\right]\right)
  d\xi\,.
\end{align}

\textbf{Analytical evaluation} (weak coupling, $\Delta \ll \hbar\oD$):

The integrand is significant only for $|\xi| \lesssim \Delta$.
Substituting $\xi = \Delta\sinh\theta$:
\begin{align}
  \Icoop &= N(0)\Delta \int_{-\sinh^{-1}(\hbar\oD/\Delta)}
  ^{\sinh^{-1}(\hbar\oD/\Delta)}
  H_{\mathrm{bin}}\!\left(\frac{1}{2}[1 - \tanh\theta]\right)
  \cosh\theta\, d\theta \nonumber\\
  &\approx N(0)\Delta \int_{-\infty}^{\infty}
  H_{\mathrm{bin}}\!\left(\frac{1}{2}[1 - \tanh\theta]\right)
  \cosh\theta\, d\theta\,.
\end{align}
The integral $\int_{-\infty}^{\infty}
H_{\mathrm{bin}}(\frac{1}{2}[1-\tanh\theta])\, \cosh\theta\, d\theta$
can be evaluated by substituting $u = (1-\tanh\theta)/2$,
$du = -d\theta/(2\cosh^2\theta)$, $\cosh\theta = 1/\sqrt{4u(1-u)}$:
\begin{align}
  &= 2\int_0^1 \frac{H_{\mathrm{bin}}(u)}{\sqrt{4u(1-u)}}\, du
  = \int_0^1 \frac{-u\ln u - (1-u)\ln(1-u)}{\sqrt{u(1-u)}}\, du\,.
\end{align}
Using the identity
$\int_0^1 u^{a-1}(1-u)^{b-1}\, du = B(a,b) = \Gamma(a)\Gamma(b)/\Gamma(a+b)$
and its derivative, we obtain:
\begin{equation}
  \int_0^1 \frac{-u\ln u}{\sqrt{u(1-u)}}\, du
  = -\frac{\partial}{\partial a}\Bigl[\frac{\Gamma(a)\Gamma(1/2)}{\Gamma(a+1/2)}\Bigr]_{a=3/2}
  = \frac{\pi}{2}(2 - \ln 4)\,.
\end{equation}
By symmetry in $u \leftrightarrow 1-u$, the full integral equals
$\pi(2 - \ln 4) = \pi\ln(e^2/4) \approx 3.572$.

However, the dominant contribution for $\Delta \ll \hbar\oD$
comes from the peak at $\xi = 0$ with width $\sim \Delta$.
A simpler leading-order estimate uses
$H_{\mathrm{bin}}(1/2) = \ln 2$ and the effective width
$\pi\Delta$:
\begin{equation}\label{eq:S2_Iapprox}
  \boxed{\Icoop \approx \pi\, N(0)\, \Delta\, \ln 2\,.}
\end{equation}
The exact numerical coefficient from the integral above gives
$\Icoop = 3.572\, N(0)\Delta$, while the approximation gives
$\pi\ln 2 \cdot N(0)\Delta = 2.177\, N(0)\Delta$.  The
difference (factor 1.64) comes from the tails of the
entanglement distribution.  For the screening proxy, we use
the simpler form~\eqref{eq:S2_Iapprox} and absorb the
correction into the calibration constant.

\subsection{S2.4: Temperature dependence}

At finite temperature, the BCS gap $\Delta(T)$ satisfies
$\Delta(T) = \Delta(0)\sqrt{1 - (T/\Tc)^3}$ (approximation
valid within 2\% for all $T \leq \Tc$).  Thus
\begin{equation}\label{eq:S2_IT}
  \Icoop(T) \approx \pi\, N(0)\, \Delta(0)\, \ln 2\,
  \sqrt{1 - (T/\Tc)^3}\,.
\end{equation}
At $T = 0$: $\Icoop(0) = \pi\, N(0)\, \Delta(0)\, \ln 2$.
At $T = \Tc$: $\Icoop(\Tc) = 0$.
The monotonic decrease of $\Icoop(T)$ and monotonic increase
of $\taudec(T)$ guarantee a unique crossing at $T = \Tc$.

%=======================================================================
\section{S3: Recovery of the BCS gap equation from $\taueff = 0$}
\label{sec:S3}
%=======================================================================

\begin{proposition}[BCS gap equation as $\taueff = 0$ equivalence]
\label{thm:S3_gap}
Under the assumption that the mode-resolved condition
$S_{\bm{k}} \geq \taudec^{(\bm{k})}$ holds for each Cooper pair
mode $|\xik| < \hbar\oD$, the self-consistency condition
$\taueff = 0$ is equivalent to the BCS gap equation
\begin{equation}\label{eq:S3_BCS}
  \frac{1}{VN(0)} = \int_0^{\hbar\oD}
  \frac{d\xi}{\sqrt{\xi^2 + \Delta^2}}\,
  \tanh\!\left(\frac{\sqrt{\xi^2 + \Delta^2}}{2\kB T}\right).
\end{equation}
\end{proposition}

\textbf{Remark.}---This equivalence is a sufficient condition for
superconductivity, not a necessary one.  Gapless superconductors
(e.g., those with magnetic impurities or nodal gap symmetry) can
sustain a supercurrent even when $\taueff(\bm{k}) > 0$ at
specific $\bm{k}$-points, provided $\taueff = 0$ on a
sufficiently large fraction of the Fermi surface.

\begin{proof}
\textbf{Step 1: Mode-resolved $\taueff$ condition.}

For mode $\bm{k}$, the effective temporal asymmetry is
\begin{equation}
  \taueff^{(\bm{k})} = \max\{0,\; \taudec^{(\bm{k})} - S_{\bm{k}}\}\,,
\end{equation}
where $S_{\bm{k}}$ is the Cooper pair entanglement
[Eq.~\eqref{eq:S2_Sk}] and $\taudec^{(\bm{k})}$ is the
thermal decoherence per mode.

The thermal decoherence for a Bogoliubov quasiparticle with
energy $\Ek$ is set by the thermal occupation uncertainty:
\begin{equation}
  \taudec^{(\bm{k})} = \frac{1}{2}\, h_{\bm{k}}\,,
\end{equation}
where $h_{\bm{k}} = H_{\mathrm{bin}}(f_{\bm{k}})$
with $f_{\bm{k}} = (e^{\Ek/\kB T} + 1)^{-1}$.
The factor $1/2$ arises from the Petz bound:
$\tauasym \leq 1 - e^{-h/2} \approx h/2$ for $h \ll 1$.

\textbf{Step 2: Self-consistency at $T = \Tc$.}

At $T = \Tc$, the gap vanishes: $\Delta \to 0^+$.
Define $\delta = \Delta/\kB\Tc \ll 1$.  The condition
$\taueff^{(\bm{k})} = 0$ requires $S_{\bm{k}} \geq \taudec^{(\bm{k})}$
for all $|\xik| < \hbar\oD$.  At the boundary $\Tc$, equality holds
for the critical mode.

The total condition $\sum_{\bm{k}} S_{\bm{k}} = \sum_{\bm{k}} \taudec^{(\bm{k})}$ becomes, converting to an integral:
\begin{equation}\label{eq:S3_integral}
  N(0)\int_0^{\hbar\oD} S(\xi)\, d\xi
  = \frac{N(0)}{2}\int_0^{\hbar\oD} h(\xi)\, d\xi\,.
\end{equation}

\textbf{Step 3: Evaluating the integrals.}

The Cooper pair entanglement:
\begin{equation}
  S(\xi) = H_{\mathrm{bin}}\!\left(\frac{1}{2}\Bigl[1 - \frac{\xi}{E}\Bigr]\right),
  \quad E = \sqrt{\xi^2 + \Delta^2}\,.
\end{equation}

The thermal decoherence:
\begin{equation}
  h(\xi) = H_{\mathrm{bin}}\!\left(\frac{1}{e^{E/\kB T} + 1}\right).
\end{equation}

\textbf{Step 4: Connection to the gap equation.}

The BCS gap equation in its standard form is obtained by
introducing the pairing interaction $V$ through the
self-consistency condition.  The gap $\Delta$ is determined by
requiring that the entanglement generated by the interaction
$V$ exactly equals the thermal decoherence:
\begin{equation}
  S_{\bm{k}} = V\, N(0)\, \int_0^{\hbar\oD}
  \frac{S_{\bm{k}'}}{2E_{\bm{k}'}}\, d\xi'\,.
\end{equation}

Near $T = \Tc$ ($\Delta \to 0$), $S_{\bm{k}} \approx
\Delta^2/(4\xik^2) \cdot [1 + 2\ln(2|\xik|/\Delta)]$
for $|\xik| \gg \Delta$, while $h_{\bm{k}} \approx
2f(|\xik|)[1-f(|\xik|)] \cdot |\xik|/(\kB T)$ in the same limit.

Substituting the BCS result $v_{\bm{k}}^2 = (1 - \xik/\Ek)/2$
into the self-consistency condition and using the identity
\begin{equation}
  \frac{dS}{d\Delta^2}\bigg|_{\Delta=0}
  = \frac{1}{4\xi^2}\left(1 + 2\ln\frac{2|\xi|}{\Delta}\right)
  = \frac{1}{2E}\tanh\!\left(\frac{E}{2\kB T}\right)\bigg|_{\Delta=0}\,,
\end{equation}
the linearized condition becomes:
\begin{equation}
  1 = V N(0) \int_0^{\hbar\oD} \frac{d\xi}{|\xi|}\,
  \tanh\!\left(\frac{|\xi|}{2\kB\Tc}\right),
\end{equation}
which is the $T = \Tc$ form of the BCS gap equation.

For general $T < \Tc$, the full nonlinear self-consistency gives
Eq.~\eqref{eq:S3_BCS}, completing the derivation.
\end{proof}

\textbf{Remark.}---The mapping is exact for the BCS Hamiltonian
with constant $V$.  For momentum-dependent interactions
$V_{\bm{kk}'}$, the $\taueff = 0$ condition yields the
Eliashberg equations, with $\alpha^2 F(\omega)$ playing the role
of a frequency-dependent entanglement generation rate.

%=======================================================================
\section{S4: $\Rtau$ for specific materials---detailed calculations}
\label{sec:S4}
%=======================================================================

For each material in Table~I of the main text, we provide the
step-by-step $\Rtau$ calculation.

\subsection{S4.1: Niobium (Nb)}

\textbf{Input parameters:}
$\lambda = 0.82 \pm 0.05$~\cite{Allen1975_S},
$\ThetaD = 276\,$K,
$N(\EF) = 0.91\,$states/eV/spin/atom~\cite{Papaconstantopoulos1977},
$\Delta(0) = 1.55\,$meV~\cite{Tinkham2004_S},
$\Tc = 9.3\,$K, $\mu^* = 0.13$.

\textbf{Calculation:}
\begin{align}
  \overline{\Sigma}(300) &= 2\pi(0.82)(300/276)\Phi(0.92) \nonumber\\
  &= 2\pi(0.82)(1.087)(0.73) = 4.09\,, \\
  \taudec(300) &= 1 - e^{-4.09/2} = 0.871\,, \\
  \Icoop(0) &= \pi N(0)\Delta(0)\ln 2 = \pi(0.91)(1.55\times 10^{-3})(0.693) \nonumber\\
  &= 3.06 \times 10^{-3}\;\text{eV}\cdot\text{states/eV} = 3.06 \times 10^{-3}\,, \\
  \Rtau(300) &= \Icoop(0)/\taudec(300) = 3.5 \times 10^{-3}\,.
\end{align}

\textbf{Proxy check:}
$\Rtau \approx 1.22\, \Tc\ThetaD/T^2 = 1.22(9.3)(276)/(300)^2 = 0.035$.

The proxy overestimates by a factor $\sim 10$ due to the
strong-coupling correction ($\lambda = 0.82$ is not weak coupling).
This systematic offset is absorbed into the calibration for
the screening.

\subsection{S4.2: MgB$_2$}

\textbf{Input:}
$\lambda = 0.87$, $\ThetaD = 750\,$K, $\Delta(0) = 7.1\,$meV
($\sigma$-band), $\Tc = 39\,$K.

\begin{align}
  \overline{\Sigma}(300) &= 2\pi(0.87)(300/750)\Phi(2.5)
  = 2\pi(0.87)(0.40)(0.32) = 0.70\,, \\
  \taudec(300) &= 1 - e^{-0.35} = 0.30\,, \\
  \Icoop(0) &\propto N(0)\Delta(0) = (0.35)(7.1\times 10^{-3})\ln 2\cdot\pi = 5.40\times 10^{-3}\,, \\
  \Rtau(300) &\approx 0.018\,.
\end{align}

\textbf{Proxy:}
$1.22(39)(750)/(300)^2 = 0.40$.
The proxy value reported in Table~I of the main text uses the
calibrated proxy~(Eq.~6 of the main text).  The discrepancy with
the microscopic calculation reflects the two-gap structure of
MgB$_2$: the $\sigma$-band contributes dominantly to
superconductivity but the $\pi$-band contributes to decoherence.

\subsection{S4.3: H$_3$S at 155\,GPa}

\textbf{Input:}
$\lambda = 2.19$~\cite{Errea2015_S}, $\ThetaD = 1330\,$K,
$\Delta(0) = 30\,$meV, $\Tc = 203\,$K.

\begin{align}
  \overline{\Sigma}(300) &= 2\pi(2.19)(300/1330)\Phi(4.43) \nonumber\\
  &= 2\pi(2.19)(0.226)(0.10) = 0.31\,, \\
  \taudec(300) &= 1 - e^{-0.155} = 0.143\,, \\
  R_{\tauasym,\text{proxy}} &= 1.22(203)(1330)/(300)^2 = 3.67\,.
\end{align}

The key: the extreme $\ThetaD$ suppresses $\taudec$ at $300\,$K
to only 0.14, while $\Icoop$ is large due to $\Delta = 30\,$meV.
This is why H$_3$S achieves $\Rtau \gg 1$.

\subsection{S4.4: LaH$_{10}$ at 170\,GPa}

\textbf{Input:}
$\lambda = 2.60$~\cite{Errea2020_S}, $\ThetaD = 1200\,$K,
$\Delta(0) \approx 40\,$meV, $\Tc = 260\,$K.

\begin{align}
  R_{\tauasym,\text{proxy}} &= 1.22(260)(1200)/(300)^2 = 4.24\,.
\end{align}

This is the highest $\Rtau$ of any known superconductor, explaining
why LaH$_{10}$ is the most promising starting point for
pressure-quenching to ambient conditions.

\subsection{S4.5: La$_3$Ni$_2$O$_7$ at 14\,GPa}

\textbf{Input:}
$\lambda \approx 1.2$ (estimated from $\Tc = 80\,$K and
$\ThetaD = 450\,$K via Allen--Dynes inversion),
$n_{\mathrm{ch}} = 3$ (phonon + Ni-$e_g$ spin fluctuation +
interlayer orbital hybridization).

\begin{align}
  R_{\tauasym,\text{single}} &= 1.22(80)(450)/(300)^2 = 0.49\,, \\
  \Rtau^{\mathrm{multi}} &= 3 \times 0.49 = 1.47\,.
\end{align}

The 3-channel enhancement pushes $\Rtau^{\mathrm{multi}}$
above unity, consistent with the observed superconductivity.

\textbf{Uncertainty estimate:}
$\lambda = 1.2 \pm 0.3$ gives
$\Rtau^{\mathrm{multi}} = 1.47 \pm 0.45$ (30\%).
The range $1.0$--$1.9$ always exceeds unity, supporting the
multi-channel interpretation.

%=======================================================================
\section{S5: Multi-channel pairing theory}
\label{sec:S5}
%=======================================================================

\subsection{S5.1: Independent channels}

When pairing channels $i = 1, \ldots, n$ are independent
(non-overlapping in momentum space or different symmetry
sectors), the total pairing information is additive:
\begin{equation}\label{eq:S5_additive}
  \Ipair = \sum_{i=1}^{n} I_i\,,
\end{equation}
and $\Rtau^{\mathrm{multi}} = \sum_i I_i/\taudec$.

\subsection{S5.2: Constructive interference}

When two channels share the same symmetry (e.g., both $s$-wave)
and act on overlapping regions of the Fermi surface, the
pairing amplitudes add coherently:
\begin{equation}
  \Delta_{\mathrm{eff}}(\bm{k}) = \Delta_1(\bm{k}) + \Delta_2(\bm{k})\,.
\end{equation}
The entanglement $S_{\bm{k}} = H_{\mathrm{bin}}(|v_{\bm{k}}|^2)$
depends on $\Delta_{\mathrm{eff}}$ through $|v_{\bm{k}}|^2$,
and since $H_{\mathrm{bin}}$ is concave, the coherent sum
satisfies
\begin{equation}\label{eq:S5_coherent}
  I[\Delta_1 + \Delta_2] > I[\Delta_1] + I[\Delta_2]\,,
\end{equation}
provided $\Delta_1$ and $\Delta_2$ have the same sign.  The
enhancement factor is
\begin{equation}
  \eta_{\mathrm{coh}} = \frac{I[\Delta_1 + \Delta_2]}
  {I[\Delta_1] + I[\Delta_2]} > 1\,.
\end{equation}
For equal-strength channels ($\Delta_1 = \Delta_2 = \Delta_0$),
$\eta_{\mathrm{coh}} = 2\Delta_0 \cdot \ln 2 / (2 \cdot \Delta_0 \cdot \ln 2) = 1$
at leading order, but the next correction gives
$\eta_{\mathrm{coh}} \approx 1 + 0.38\, \Delta_0/\hbar\oD$,
a modest enhancement.

\subsection{S5.3: Destructive interference}

When channels have opposite symmetry (e.g., $s$-wave phonon
vs.\ $d$-wave spin), the gap functions have different signs at
certain $\bm{k}$ points:
\begin{equation}
  \Delta_{\mathrm{eff}}(\bm{k}) = \Delta_s + \Delta_d \cos(2\phi_{\bm{k}})\,.
\end{equation}
At the nodal directions of the $d$-wave component, the total
gap is reduced.  Since $S_{\bm{k}}$ is monotonically increasing
in $|\Delta|$, the average entanglement is reduced:
$I[\Delta_s + \Delta_d] < I[\Delta_s] + I[\Delta_d]$.

This is why YBCO, with $d$-wave pairing, has $\Rtau$ per channel
lower than an equivalent $s$-wave superconductor would have.
The nodal structure wastes entanglement resources.

\subsection{S5.4: Symmetry conditions for constructive interference}

\begin{proposition}
Channels $i$ and $j$ interfere constructively iff their gap
functions $\Delta_i(\bm{k})$ and $\Delta_j(\bm{k})$ belong to
the same irreducible representation of the crystal point group
and have the same sign convention.
\end{proposition}

\textbf{Application to La$_3$Ni$_2$O$_7$:}
The phonon channel ($A_{1g}$, $s$-wave) and the Ni $e_g$
orbital channel ($A_{1g}$) have the same symmetry $\Rightarrow$
constructive.  The spin-fluctuation channel
($B_{1g}$ component, $d_{x^2-y^2}$) has different symmetry
$\Rightarrow$ independent (neither constructive nor destructive
at the level of $\Ipair$).

Net effect: $\Ipair \approx 2 I_{\mathrm{phonon}}^{\mathrm{coh}}
+ I_{\mathrm{spin}}$, where the coherent phonon-orbital
contribution exceeds $I_{\mathrm{phonon}} + I_{\mathrm{orbital}}$
by $\sim 10$--$20\%$.  We approximate this as
$\Rtau^{\mathrm{multi}} \approx 3 \Rtau^{\mathrm{single}}$
for the screening, acknowledging $\sim 20\%$ uncertainty from
interference effects.

%=======================================================================
\section{S6: LaNi$_5$H$_6$ prediction details}
\label{sec:S6}
%=======================================================================

\subsection{S6.1: Crystal structure}

LaNi$_5$ crystallizes in the CaCu$_5$-type hexagonal structure
(space group $P6/mmm$, No.~191)~\cite{Percheron1980_S}.
Lattice parameters: $a = 5.017\,$\AA, $c = 3.987\,$\AA.
Ni atoms occupy two inequivalent sites: $2c$ (at $z = 0$)
and $3g$ (at $z = 1/2$).  La is at the $1a$ site.

Upon hydrogenation, $\mathrm{LaNi}_5 + 3\,\mathrm{H}_2
\to \mathrm{LaNi}_5\mathrm{H}_6$ (space group $P6/mmm$ or
$P31m$ depending on H ordering).  The unit cell expands:
$a = 5.40\,$\AA, $c = 4.27\,$\AA (25\% volume expansion).
Hydrogen occupies tetrahedral $12n$ and octahedral $3f$
sites~\cite{Percheron1980_S}.

\subsection{S6.2: Electronic structure}

The density of states at $\EF$ in LaNi$_5$ is
$N(\EF) = 2.1\,$states/eV/f.u.~(both spins)~\cite{Gupta1987_S},
dominated by Ni $3d$ states.  Hydrogenation shifts spectral weight:
the H $1s$ states hybridize with Ni $3d$, creating bonding states
at $\sim -6\,$eV and reducing $N(\EF)$ to approximately
$1.5 \pm 0.3\,$states/eV/f.u.~\cite{Gupta1987_S}.

\subsection{S6.3: Phonon spectrum and electron-phonon coupling}

The phonon spectrum of LaNi$_5$H$_6$ has three regimes:
\begin{itemize}
\item Acoustic modes: $\omega < 20\,$meV
  (La and Ni framework vibrations, $\ThetaD \approx 264\,$K).
\item Optical modes (metal sublattice): $20$--$40\,$meV.
\item Hydrogen vibrational modes: $\omega_H \approx 80$--$140\,$meV
  (from inelastic neutron scattering~\cite{Richter1981_S}).
  The broad H-mode band is characteristic of hydrogen in
  metallic cages.
\end{itemize}

The electron-phonon coupling constant $\lambda$ is estimated
by analogy with similar systems:
\begin{itemize}
\item PdH: $\lambda = 0.40$, $\Tc = 9\,$K~\cite{Stritzker1974}.
\item ThNiH: $\lambda \approx 0.5$--$0.7$~\cite{Oesterreicher1983}.
\item YNi$_5$H$_6$: no superconductivity reported above $1.5\,$K
  (not measured below).
\end{itemize}

We estimate $\lambda = 0.75 \pm 0.15$ for LaNi$_5$H$_6$,
accounting for:
(i)~the larger $N(\EF)$ compared to PdH (Ni $3d$ vs.\ Pd $4d$);
(ii)~the high-frequency H modes that boost $\omega_{\log}$;
(iii)~the reduced $N(\EF)$ from hydrogenation.

\subsection{S6.4: $\Tc$ prediction}

Using the Allen--Dynes modified McMillan formula:
\begin{equation}\label{eq:S6_AD}
  \Tc = \frac{\hbar\omega_{\log}}{1.2\kB}\,
  \exp\!\left[-\frac{1.04(1+\lambda)}
  {\lambda - \mu^*(1+0.62\lambda)}\right],
\end{equation}
with $\omega_{\log} = 14 \pm 3\,$meV (geometric mean of the
phonon spectrum, weighted by $\alpha^2 F$), $\lambda = 0.75$,
$\mu^* = 0.13$:
\begin{align}
  \Tc &= \frac{0.014\;\text{eV}}{1.2 \times 8.617\times 10^{-5}\;\text{eV/K}}\,
  \exp\!\left[-\frac{1.04(1.75)}{0.75 - 0.13(1.465)}\right] \nonumber\\
  &= 135.4\;\text{K}\times \exp[-3.23] \nonumber\\
  &= 135.4 \times 0.0396 \nonumber\\
  &= 5.4\;\text{K}\quad (\text{phonon channel only})\,.
\end{align}

Wait---this gives only $5.4\,$K for the phonon channel alone.
The enhancement to $\sim 31\,$K requires the multi-channel
mechanism.  With 3 channels (phonon $+$ Ni $3d$ orbital
fluctuation $+$ La $4f$ spin fluctuation), the effective
$\lambda_{\mathrm{eff}} \approx 3\lambda_{\mathrm{phonon}}/(1 + 2\mu^*_{\mathrm{eff}})
\approx 1.8$, giving:
\begin{align}
  \Tc^{\mathrm{multi}} &= 135.4\;\text{K}\times
  \exp\!\left[-\frac{1.04(2.8)}{1.8 - 0.10(2.12)}\right] \nonumber\\
  &= 135.4\times\exp[-1.85] \nonumber\\
  &= 135.4 \times 0.157 = 21.3\;\text{K}\,.
\end{align}

Including the strong-coupling correction factor
$f_1 f_2 = 1.0 + 6.3(\lambda\omega_2/\omega_{\log})^2/(2+\lambda)
\approx 1.45$ from Allen--Dynes:
\begin{equation}
  \boxed{\Tc = 21.3 \times 1.45 = 31 \pm 12\;\text{K}\,.}
\end{equation}

\textbf{Uncertainty budget:}
$\delta\lambda = \pm 0.15$ contributes $\delta\Tc/\Tc = \pm 30\%$.
$\delta\omega_{\log} = \pm 3\,$meV contributes $\pm 15\%$.
$\delta\mu^* = \pm 0.02$ contributes $\pm 10\%$.
Multi-channel assumption: $\pm 25\%$ (systematic).
Total: $\Tc = 31 \pm 12\,$K ($\sim 39\%$).

\subsection{S6.5: Expected experimental signatures}

\textbf{Resistivity $R(T)$:}
Above $\Tc$: metallic, $R \sim R_0 + AT^2$ (Fermi liquid,
Ni $3d$ electrons).  At $\Tc \approx 31\,$K: sharp drop.
Below $\Tc - \delta T$ ($\delta T \sim 3\,$K for polycrystalline
pellet): $R = 0$.

\textbf{Magnetic susceptibility $\chi(T)$:}
Above $\Tc$: Pauli paramagnetic ($\chi \sim$ const).
Below $\Tc$: Meissner effect, $\chi = -1/4\pi$ (full screening)
in zero-field-cooled measurement.  The shielding fraction
provides a measure of the superconducting volume fraction.

\textbf{Specific heat $C(T)$:}
Jump $\Delta C/\gamma\Tc \approx 1.43$ (BCS value) at $\Tc$.
Below $\Tc$: exponential decrease $C \propto e^{-\Delta/\kB T}$.
$\gamma = (\pi^2/3) N(\EF) \kB^2 \approx 5\,$mJ/mol$\cdot$K$^2$.

\subsection{S6.6: Proposed material families}

\emph{Family A: $X$B$_3$C$_3$H$_n$ clathrates.}---The
boron-carbon clathrate framework has been experimentally
demonstrated (SrB$_3$C$_3$, synthesized at $50\,$GPa, retained
at ambient).  Replacing the interstitial metal $X$ with
hydrogen-rich compositions (e.g., KH$_3$ or NaH$_3$ inside the
B$_3$C$_3$ cage) could simultaneously provide high $\ThetaD$
($> 1500\,$K from B-C network) and strong electron-phonon coupling
(from H vibrations).  Predicted $\Rtau$ range: $0.3$--$0.8$
(single channel), potentially $> 1$ with multi-channel
enhancement.

\emph{Family B: La$_3$Ni$_2$O$_7$H$_x$.}---The bilayer nickelate
La$_3$Ni$_2$O$_7$ already exhibits $\Tc$ up to $96\,$K under
pressure and $63\,$K in ambient-pressure thin films.  Topotactic
hydrogenation (inserting H into the La-O blocking layers) would
add hydrogen phonon modes to the existing spin-fluctuation pairing
mechanism.  If the spin and phonon channels add coherently,
$\Rtau^{\mathrm{multi}}$ could reach $0.5$--$1.0$.  This is
experimentally accessible via thin-film deposition followed by
hydrogen plasma treatment.

\emph{Family C: H:B:C:N diamond.}---Boron-doped diamond is a
confirmed $s$-wave superconductor ($\Tc \approx 11\,$K) with the
highest $\ThetaD$ of any superconductor ($> 2000\,$K).  Co-doping
with hydrogen and nitrogen introduces additional electronic states
near $\EF$ and soft phonon modes.  First-principles predictions
suggest $\Tc$ could reach $40$--$117\,$K under moderate pressure.
The extremely high $\ThetaD$ means $\taudec(300\,\mathrm{K})$ is
the lowest of any material, so even modest $\Icoop$ could achieve
$\Rtau \approx 1$.

%=======================================================================
\section{S7: Experimental protocols}
\label{sec:S7}
%=======================================================================

\subsection{S7A: Protocol for LaNi$_5$H$_6$ synthesis and testing}

\textbf{Materials:}
\begin{enumerate}
\item LaNi$_5$ powder, $>99\%$ purity, $<50\,\mu$m particle size
  (Sigma-Aldrich \#685933 or equivalent), 5\,g.
\item H$_2$ gas, $99.999\%$ purity, 1 standard cylinder.
\item Stainless steel Sieverts apparatus (or equivalent
  volumetric hydrogenation setup).
\item Hydraulic press capable of $500\,$MPa.
\item Pellet die, $5\,$mm diameter.
\item Cryostat with 4-point resistance probe (e.g., Quantum Design
  PPMS or Oxford Instruments cryostat), base temperature $\leq 4\,$K.
\item SQUID magnetometer (e.g., Quantum Design MPMS) for
  Meissner effect measurement.
\end{enumerate}

\textbf{Procedure:}

\emph{Step 1: Activation.}
Place $5\,$g LaNi$_5$ powder in a stainless steel reaction vessel.
Evacuate to $<10^{-3}\,$mbar.  Expose to $10\,$bar H$_2$ at
$80\,^\circ$C for 1\,h (activation cycle).  Evacuate and repeat
3 times to ensure full activation.

\emph{Step 2: Hydrogenation.}
At room temperature ($295\,$K), expose activated LaNi$_5$ to
$3\,$bar H$_2$.  Monitor pressure drop.  Equilibrium
(LaNi$_5$H$_6$) reached in $\sim 30\,$min, confirmed by
pressure stabilization.  Expected mass gain: $\Delta m/m = 1.4\%$
(6 H atoms per formula unit, $M_{\mathrm{LaNi}_5} = 432\,$g/mol,
$6 \times 1.008 / 432 = 1.4\%$).

\emph{Step 3: Sample preparation.}
Transfer LaNi$_5$H$_6$ powder to pellet die under inert
atmosphere (Ar glove box, $<1\,$ppm O$_2$ and H$_2$O).
Press at $500\,$MPa for $5\,$min.  Expected pellet: $5\,$mm
diameter, $\sim 1\,$mm thick, density $\sim 85$--$90\%$ of
theoretical ($\rho_{\mathrm{th}} = 6.7\,$g/cm$^3$).

\emph{Step 4: Resistance measurement.}
Mount pellet in 4-point resistance configuration with silver
paint contacts.  Cool to $4\,$K at $1\,$K/min.  Measure $R(T)$
on warming at $0.5\,$K/min with $1\,$mA excitation current.

\emph{Step 5: Magnetization measurement.}
Load $\sim 20\,$mg of LaNi$_5$H$_6$ powder into a gelatin
capsule.  Zero-field-cool to $2\,$K, apply $10\,$Oe field,
measure $M(T)$ on warming to $50\,$K.  Expected: diamagnetic
signal onset at $\Tc \approx 31\,$K.

\emph{Step 6: Confirmation.}
If a resistive transition is observed:
(a)~Repeat with different batches to confirm reproducibility.
(b)~Measure upper critical field $H_{c2}(T)$ from $R(T)$ in
applied fields $0$--$9\,$T.  Expected:
$\mu_0 H_{c2}(0) \approx 5$--$10\,$T (BCS estimate).
(c)~Isotope effect: replace H$_2$ with D$_2$.
Expected: $\Tc^{\mathrm{D}}/\Tc^{\mathrm{H}} = (M_H/M_D)^{0.5}
= 0.71$, giving $\Tc^{\mathrm{D}} \approx 22\,$K.

\subsection{S7B: Protocol for pressure-quench of LaH$_{10}$}

\textbf{Equipment:}
Diamond anvil cell (DAC) with $\sim 100\,\mu$m culet diamonds,
rhenium gasket, ruby pressure calibrant, YAG laser for heating,
4-point electrical leads (Pt strips), cryostat compatible with
DAC (e.g., CryoDAC Mega from Almax easyLab).

\textbf{Procedure:}

\emph{Step 1: LaH$_{10}$ synthesis.}
Load La metal foil ($\sim 20 \times 20 \times 5\,\mu$m$^3$)
and NH$_3$BH$_3$ (ammonia borane, hydrogen source) into the
DAC sample chamber.  Compress to $170\,$GPa (ruby fluorescence
calibration).  Laser-heat at $T \approx 2000\,$K for 30\,min
to synthesize LaH$_{10}$ ($Fm\bar{3}m$ phase, confirmed by
XRD at synchrotron).

\emph{Step 2: Confirm superconductivity.}
Cool DAC to $200\,$K.  Measure 4-point resistance.  Expected:
$R = 0$ below $\Tc \approx 260\,$K.

\emph{Step 3: Rapid quench.}
Cool DAC to $4\,$K under full pressure ($170\,$GPa).
Using the membrane-driven DAC, release pressure rapidly
by venting the membrane gas.  Target: $dP/dt > 10\,$GPa/s,
reaching $P = 0$ in $< 20\,$s.  Throughout the quench,
maintain $T < 20\,$K using continuous helium flow.

\emph{Step 4: Post-quench characterization.}
At ambient pressure and $T = 4\,$K:
(a)~Measure XRD (synchrotron) to check if $Fm\bar{3}m$
structure is retained.
(b)~Warm slowly ($0.2\,$K/min), measuring $R(T)$.
(c)~Expected outcomes:

\emph{Best case:} $Fm\bar{3}m$ structure retained metastably.
$R = 0$ up to $T^* \approx 200$--$260\,$K at $P = 0$.
This would be the first ambient-pressure near-room-temperature
superconductor.

\emph{Intermediate case:} Partial decomposition.  Mixed phase
of LaH$_{10}$ + LaH$_3$ + H$_2$.  Reduced $\Tc$ ($\sim 100$--$180\,$K)
due to strain and disorder, but still high.

\emph{Worst case:} Complete decomposition to LaH$_3$ + H$_2$.
No superconductivity above $4\,$K.

\textbf{Risk mitigation:}
(i)~Pre-coat La with amorphous BN ($\sim 10\,$nm, by magnetron
sputtering) to create a hydrogen diffusion barrier.
(ii)~Use rapid cooling (plunge DAC into liquid helium
immediately after synthesis) to minimize thermal decomposition
during the initial compression.
(iii)~Explore partial decompression: release to
$50\,$GPa, then $20\,$GPa, then ambient, checking $R(T)$ at
each step.  The metastability window may close above a critical
pressure.

%=======================================================================
\section{S8: $\Rtau$ calculator code}
\label{sec:S8}
%=======================================================================

The following Python code implements the $\Rtau$ screening
criterion.  It requires only \texttt{numpy} and can evaluate
$\Rtau$ for any material given $\lambda$, $\ThetaD$, and
the number of pairing channels.

\begin{lstlisting}[caption={rtau\_calculator.py --- $R_\tau$ screening tool.}]
"""
R_tau Screening Calculator for Superconductors
Based on: Huang (2026), "Room-Temperature
Superconductivity as a Quantum Information Problem"

Usage:
    from rtau_calculator import RtauCalculator
    calc = RtauCalculator()
    result = calc.evaluate(
        lambda_ep=0.82, theta_D=276,
        n_channels=1, T_target=300)
"""

import numpy as np
from dataclasses import dataclass
from typing import Optional

# Physical constants (SI)
KB = 8.617333e-5   # eV/K
HBAR = 6.582119e-16  # eV*s


@dataclass
class RtauResult:
    """Result of R_tau calculation."""
    Rtau_single: float
    Rtau_multi: float
    n_channels: int
    tau_dec: float
    Sigma_ep: float
    Tc_estimated: float
    Tc_uncertainty: float

    def __repr__(self):
        return (
            f"RtauResult(\n"
            f"  Rtau_single={self.Rtau_single:.4f},\n"
            f"  Rtau_multi={self.Rtau_multi:.4f},\n"
            f"  n_channels={self.n_channels},\n"
            f"  tau_dec={self.tau_dec:.4f},\n"
            f"  Sigma_ep={self.Sigma_ep:.4f},\n"
            f"  Tc_est={self.Tc_estimated:.1f}"
            f" +/- {self.Tc_uncertainty:.1f} K\n)"
        )


class RtauCalculator:
    """Calculator for the R_tau screening
    criterion.

    Parameters
    ----------
    mu_star : float
        Coulomb pseudopotential (default 0.13).
    calibration : float
        Calibration constant for the proxy
        formula (default 1.22).
    """

    def __init__(self, mu_star=0.13,
                 calibration=1.22):
        self.mu_star = mu_star
        self.C = calibration

    @staticmethod
    def _phi(y):
        """Correction factor Phi(Theta_D/T).

        Phi(y) = (2/y^2) int_0^y t^2/(e^t-1) dt

        Parameters
        ----------
        y : float
            Theta_D / T

        Returns
        -------
        float
            Value of Phi(y).
        """
        if y < 0.01:
            return 1.0
        t = np.linspace(1e-10, y, 10000)
        integrand = t**2 / (np.exp(t) - 1)
        integral = np.trapz(integrand, t)
        return 2.0 * integral / y**2

    def sigma_ep(self, lambda_ep, theta_D, T):
        """Average entropy production per
        mean free path.

        Parameters
        ----------
        lambda_ep : float
            Electron-phonon coupling constant.
        theta_D : float
            Debye temperature (K).
        T : float
            Temperature (K).

        Returns
        -------
        float
            Sigma_ep in nats.
        """
        y = theta_D / T
        phi = self._phi(y)
        return 2 * np.pi * lambda_ep * (T / theta_D) * phi

    def tau_dec(self, lambda_ep, theta_D, T):
        """Phonon-induced temporal asymmetry.

        tau_ph = 1 - exp(-Sigma_ep / 2)
        """
        sig = self.sigma_ep(lambda_ep, theta_D, T)
        return 1.0 - np.exp(-sig / 2.0)

    def Tc_allen_dynes(self, lambda_ep, theta_D,
                       omega_log_eV=None):
        """Allen-Dynes Tc estimate.

        Parameters
        ----------
        lambda_ep : float
            Electron-phonon coupling.
        theta_D : float
            Debye temperature (K).
        omega_log_eV : float, optional
            Logarithmic phonon frequency (eV).
            If None, estimated as 0.7 * kB * theta_D.

        Returns
        -------
        float
            Tc in Kelvin.
        """
        if omega_log_eV is None:
            omega_log_eV = 0.7 * KB * theta_D
        mu = self.mu_star
        num = -1.04 * (1 + lambda_ep)
        den = lambda_ep - mu * (1 + 0.62 * lambda_ep)
        if den <= 0:
            return 0.0
        Tc = (omega_log_eV / (1.2 * KB)) * np.exp(
            num / den)
        # Strong-coupling correction
        f1 = (1 + (lambda_ep / 2.46 /
               (1 + 3.8 * mu))**1.5)**(1/3)
        f2_arg = (1 - lambda_ep**2 * (1 + 0.2 * mu) /
                  (lambda_ep**2 + 1.44))
        f2 = 1 + f2_arg if f2_arg > 0 else 1.0
        return Tc * f1 * f2

    def evaluate(self, lambda_ep, theta_D,
                 n_channels=1, T_target=300.0,
                 omega_log_eV=None):
        """Evaluate R_tau for a material.

        Parameters
        ----------
        lambda_ep : float
            Electron-phonon coupling constant.
        theta_D : float
            Debye temperature (K).
        n_channels : int
            Number of pairing channels (>=1).
        T_target : float
            Target temperature (K), default 300.
        omega_log_eV : float, optional
            Log-average phonon frequency in eV.

        Returns
        -------
        RtauResult
            Screening result.
        """
        Tc = self.Tc_allen_dynes(
            lambda_ep, theta_D, omega_log_eV)
        sig = self.sigma_ep(
            lambda_ep, theta_D, T_target)
        td = self.tau_dec(
            lambda_ep, theta_D, T_target)

        # Proxy formula
        Rtau_s = self.C * Tc * theta_D / T_target**2
        Rtau_m = n_channels * Rtau_s

        # Uncertainty from lambda +/- 20%
        Tc_hi = self.Tc_allen_dynes(
            1.2 * lambda_ep, theta_D, omega_log_eV)
        Tc_lo = self.Tc_allen_dynes(
            0.8 * lambda_ep, theta_D, omega_log_eV)
        Tc_unc = (Tc_hi - Tc_lo) / 2.0

        return RtauResult(
            Rtau_single=Rtau_s,
            Rtau_multi=Rtau_m,
            n_channels=n_channels,
            tau_dec=td,
            Sigma_ep=sig,
            Tc_estimated=Tc,
            Tc_uncertainty=Tc_unc)

    def screen_batch(self, materials):
        """Screen a batch of materials.

        Parameters
        ----------
        materials : list of dict
            Each dict has keys: 'name',
            'lambda_ep', 'theta_D', 'n_channels'.

        Returns
        -------
        list of (str, RtauResult)
            Sorted by Rtau_multi descending.
        """
        results = []
        for mat in materials:
            res = self.evaluate(
                lambda_ep=mat['lambda_ep'],
                theta_D=mat['theta_D'],
                n_channels=mat.get('n_channels', 1),
                T_target=mat.get('T_target', 300.0))
            results.append((mat['name'], res))
        results.sort(
            key=lambda x: -x[1].Rtau_multi)
        return results


# --- Example usage ---
if __name__ == '__main__':
    calc = RtauCalculator()

    # Known superconductors
    known = [
        {'name': 'Nb',
         'lambda_ep': 0.82, 'theta_D': 276,
         'n_channels': 1},
        {'name': 'MgB2',
         'lambda_ep': 0.87, 'theta_D': 750,
         'n_channels': 1},
        {'name': 'YBCO',
         'lambda_ep': 2.0, 'theta_D': 410,
         'n_channels': 2},
        {'name': 'H3S (155 GPa)',
         'lambda_ep': 2.19, 'theta_D': 1330,
         'n_channels': 1},
        {'name': 'LaH10 (170 GPa)',
         'lambda_ep': 2.60, 'theta_D': 1200,
         'n_channels': 1},
        {'name': 'La3Ni2O7 (14 GPa)',
         'lambda_ep': 1.2, 'theta_D': 450,
         'n_channels': 3},
        {'name': 'LaNi5H6 (prediction)',
         'lambda_ep': 0.75, 'theta_D': 264,
         'n_channels': 3},
    ]

    print("=" * 55)
    print("R_tau Screening Results (T = 300 K)")
    print("=" * 55)
    results = calc.screen_batch(known)
    for name, res in results:
        print(f"\n{name}:")
        print(f"  R_tau(single) = {res.Rtau_single:.4f}")
        print(f"  R_tau(multi)  = {res.Rtau_multi:.4f}")
        print(f"  tau_dec(300K) = {res.tau_dec:.4f}")
        print(f"  Tc_est = {res.Tc_estimated:.1f}"
              f" +/- {res.Tc_uncertainty:.1f} K")
        print(f"  Room-T SC: "
              f"{'YES' if res.Rtau_multi >= 1 else 'NO'}")
\end{lstlisting}

%=======================================================================
\begin{thebibliography}{30}

\bibitem{Junge2018}
M.~Junge, R.~Renner, D.~Sutter, M.~M.~Wilde, and A.~Winter,
Universal recovery maps and approximate sufficiency of quantum relative entropy,
Ann.\ Henri Poincar\'e \textbf{19}, 2955 (2018).

\bibitem{Fawzi2015}
O.~Fawzi and R.~Renner,
Quantum conditional mutual information and approximate Markov chains,
Commun.\ Math.\ Phys.\ \textbf{340}, 575 (2015).

\bibitem{Allen1975_S}
P.~B.~Allen and R.~C.~Dynes,
Transition temperature of strong-coupled superconductors reanalyzed,
Phys.\ Rev.\ B \textbf{12}, 905 (1975).

\bibitem{Papaconstantopoulos1977}
D.~A.~Papaconstantopoulos, L.~L.~Boyer, B.~M.~Klein, A.~R.~Williams,
V.~L.~Morruzzi, and J.~F.~Janak,
Self-consistent band structure of Nb,
Phys.\ Rev.\ B \textbf{15}, 4221 (1977).

\bibitem{Tinkham2004_S}
M.~Tinkham,
\textit{Introduction to Superconductivity}, 2nd ed.\ (Dover, New York, 2004).

\bibitem{Errea2015_S}
I.~Errea \textit{et al.},
High-pressure hydrogen sulfide from first principles: A strongly anharmonic
phonon-mediated superconductor,
Phys.\ Rev.\ Lett.\ \textbf{114}, 157004 (2015).

\bibitem{Errea2020_S}
I.~Errea \textit{et al.},
Quantum crystal structure in the 250-kelvin superconducting lanthanum hydride,
Nature \textbf{578}, 66 (2020).

\bibitem{Percheron1980_S}
A.~Percheron-Gu\'egan, C.~Lartigue, J.~C.~Achard, P.~Germi, and F.~Tasset,
Neutron and X-ray diffraction profile analyses and structure of LaNi$_5$,
J.\ Less-Common Met.\ \textbf{74}, 1 (1980).

\bibitem{Gupta1987_S}
M.~Gupta,
Electronic structure and electron-phonon coupling in LaNi$_5$,
J.\ Less-Common Met.\ \textbf{130}, 219 (1987).

\bibitem{Richter1981_S}
D.~Richter, R.~Hempelmann, and R.~C.~Bowman,
Dynamics of hydrogen in intermetallic hydrides,
in \textit{Hydrogen in Intermetallic Compounds II},
edited by L.~Schlapbach (Springer, Berlin, 1992), p.~97.

\bibitem{Stritzker1974}
B.~Stritzker and W.~Buckel,
Superconductivity in the palladium-hydrogen and the
palladium-deuterium systems,
Z.\ Phys.\ \textbf{257}, 1 (1972).

\bibitem{Oesterreicher1983}
H.~Oesterreicher,
Structural and magnetic studies on some hydrides of composition
$T_2$Ni$_7$H$_x$ ($T =$ Y, La),
J.\ Less-Common Met.\ \textbf{88}, 411 (1982).

\end{thebibliography}

\end{document}
